\section{Conclusion}
\label{sec:conclusion}

This paper presents a comprehensive evaluation of the multi-agent collaboration framework, demonstrating its effectiveness in enabling coordinated problem-solving across a variety of enterprise domains. The results highlight the system's ability to leverage specialized agents to tackle complex tasks, achieving impressive goal success rates of up to 90\% in the evaluated scenarios. A key aspect of MAC is the focus on optimizing the inter-agent communication mechanisms. The introduction of the payload referencing feature, for example, was shown to provide significant benefits, particularly in code-heavy tasks, by allowing the supervisor agent to more efficiently share and reference large content blocks. This optimization  has led to remarkable improvements in both the overall goal success rate and the communication efficiency of the system.

Furthermore, the evaluation framework that is employed in this study, which combines assertion-based benchmarking with automated LLM-based judgements, has proven to be a reliable and scalable approach for assessing the performance of multi-agent collaborative systems. During development, we observe high agreement rates on goal success between human and automated evaluation. This validates our evaluation framework, which can enable faster prototyping and development of MAC.

One key focus for future work will be on further reducing the latency observed in more complex scenarios, such as those in the Software Development domain. Exploring additional optimizations to the inter-agent coordination mechanisms may help to address this challenge and ensure the system can operate efficiently, even in time-sensitive applications. Additionally, expanding the detection and handling of different forms of static long-form payloads could lead to additional performance gains. Investigating automatic prompt optimization techniques may also prove valuable in enhancing the collaboration capabilities of the individual agents.

Finally, automating the dataset curation process for the benchmarking framework could enable more scalable and iterative development of the multi-agent system, allowing for faster prototyping and deployment of improvements. By building on the solid foundations laid out in this technical report, the future work in these areas promises to further strengthen MAC and its ability to tackle increasingly complex, real-world challenges.
